\section{Dynamic spin polarisation, topology rules, Ovchinnikov and Lieb}\label{sec:spinpol}

\begin{reviewclaim}
``\dots instantaneous correlation effects, in particular the dynamic spin polarization (Scheme 1). A
damped spin wave so formed generates ferromagnetic coupling between the spins separated by an odd number
of sites, and antiferromagnetic between those separated by an even number of sites, thus yielding the
topology rules of spin coupling. Within the spin-Hamiltonian VB theory, Ovchinnikov's qualitative rule
dictates that the most spin-alternant determinant is expected to have the lowest energy and the highest
contribution.'' (p.~443)
\end{reviewclaim}

\subsection{Spin polarisation is a correlation effect}
\begin{enumerate}
  \item In Hückel (independent-electron) theory the unpaired electron of an odd alternant radical sits in
        the NBMO. By \cref{sec:huckel} the NBMO vanishes on the smaller (unstarred) set. The spin density on
        unstarred atoms is therefore \emph{exactly zero}.
  \item Add on-site repulsion $U\sum_i\hat n_{i\uparrow}\hat n_{i\downarrow}$ (Hubbard model). An $\uparrow$ electron
        in the NBMO repels $\downarrow$ density on the starred atoms more than $\uparrow$ density. Paired
        electrons therefore rearrange so that $\downarrow$ density accumulates on the neighbouring (unstarred)
        atoms: a negative spin density appears between positive ones. This is Scheme~1 of the review.
\end{enumerate}
\begin{verifybox}
T9: pentadienyl radical (5-site chain, 5 electrons, $S_z=\tfrac12$), exact diagonalisation.
$U=0$: $\rho_s=(0.333,0,0.333,0,0.333)$. $U=3t$: $\rho_s=(0.451,-0.182,0.462,-0.182,0.451)$, sum $=1$.
The alternation appears only with $U\neq0$. \textbf{PASS}.
\end{verifybox}

\subsection{Odd/even rule from the bipartite structure}
\begin{enumerate}
  \item From \cref{sec:kinex}, for $U\gg t$ each bond carries an antiferromagnetic coupling
        $J_{\text{AF}}=4t^2/U$. At half filling the low-energy physics is the Heisenberg model
        $\hat H=J_{\text{AF}}\sum_{\langle ij\rangle}\hat{\mathbf S}_i\cdot\hat{\mathbf S}_j$.
  \item On a bipartite graph every bond joins $\mathcal A$ to $\mathcal B$. The energy is lowest when neighbours
        are antiparallel: spins on $\mathcal A$ point one way and spins on $\mathcal B$ the other (the Néel pattern).
  \item Two sites are in the same set iff the path between them has an even number of bonds, i.e.\ an
        \emph{odd} number of intervening atoms. Hence:
        same set $\Rightarrow$ parallel (ferromagnetic); different sets $\Rightarrow$ antiparallel.
  \item Example: in $m$-xylylene the radical CH$_2$ carbons are joined via 3 intervening ring atoms (odd), so
        they couple ferromagnetically (triplet). In $p$-xylylene there are 4 (even), so the coupling is
        antiferromagnetic: a singlet (actually Kekulé, closed shell). This is the review's rule.
\end{enumerate}

\subsection{Ovchinnikov's rule and Lieb's theorem}
\begin{enumerate}
  \item Counting the Néel configuration: $n_A$ spins up and $n_B$ spins down gives $S=\tfrac12|n_A-n_B|$.
        This is \textbf{Ovchinnikov's rule} for the ground-state spin of an alternant hydrocarbon,
        $S=\tfrac12|n^*-n^\circ|$.
  \item The rule was conjectured from VB/Heisenberg arguments. Lieb (1989) proved it rigorously for the
        repulsive Hubbard model at half filling on any bipartite graph: \emph{the ground state is unique and
        has $S=\tfrac12\bigl||\mathcal B|-|\mathcal A|\bigr|$}, for every $U>0$ [verified: \emph{Phys.\ Rev.\ Lett.}\ \textbf{62}, 1201 (1989)].
  \item Consequences: CBD and TME ($n_A=n_B$) have singlet ground states; TMM (1/3) and MX (3/5) are triplets;
        TriMB (3/6) is a quartet; 1,2,4,5-tetramethylenebenzene (5/5) is a singlet. These are the review's
        Figure~3 assignments.
\end{enumerate}
\begin{verifybox}
T9: Hubbard model ($t=1$, $U=3$), exact diagonalisation in every $S_z$ sector. The ground-state $S$ equals
$\tfrac12|n_A-n_B|$ for all 7 graphs: CBD 0, TMM 1, TME 0, MX 1, TriMB $\tfrac32$, TetraMB 0, $p$-xylylene 0. \textbf{PASS}.
\end{verifybox}

\subsection{Disjoint diradicals: the singlet is selected by spin polarisation (Kollmar--Staemmler)}
For disjoint NBMOs (CBD, TME) the first-order gap $2K_{ab}$ vanishes (\cref{sec:huckel}). The degeneracy is
lifted only at second order, through the spin-polarised configurations above. In D$_{4h}$ CBD, ``the
bonding $\pi$-electrons can partially localize on the same pair of atoms as the same-spin unpaired
electrons; however, in the triplet state such advantageous spin polarization is not possible'' (review,
p.~443). The numerical signature is a gap $\propto U^2$, with the singlet lower.
\begin{verifybox}
T9b: gap $E_T-E_S>0$ (singlet lower), scaling as $U^{1.97}$ (CBD) and $U^{1.95}$ (TME) for $U\to0$;
cf.\ linear scaling with the triplet lower for the nondisjoint TMM and MX. \textbf{PASS}.
\end{verifybox}
\begin{sourcebox}
Kollmar \& Staemmler, \emph{Theor.\ Chim.\ Acta} \textbf{48}, 223 (1978), doi:10.1007/BF00549021 (not accessed).
Ovchinnikov, \emph{Theor.\ Chim.\ Acta} \textbf{47}, 297 (1978), doi:10.1007/BF00549259: $S=(N_A-N_B)/2$ as quoted in
arXiv:1906.08544. Lieb 1989: statement verified.
\end{sourcebox}
\begin{caveatbox}
Ovchinnikov's rule and Lieb's theorem apply to \emph{alternant} (bipartite) $\pi$ systems described by
Hubbard/Heisenberg-type models. The review's own counterexamples (pyridinium analogue V, Figure~6;
$\sigma$--$\pi$ diradicals such as the didehydrotoluenes) lie outside these assumptions. This is consistent
with its warning that such rules ``inevitably'' have ``limited scope''.
\end{caveatbox}
